% ------------------------------------------------------------------
%  Capitolo 14 — Il metodo in una settimana
%  Parte III
%  Ricerca di riferimento: ricerca 01 §17.2; 04 §6; 03 §4.4
%  Voci bibliografiche principali: dunlosky2013, donoghue2021, carpenter2022
%  Epigrafe: ✔ verificata sul testo (vedi ricerca 03 e registro, ricerca 00)
%  Scheda del capitolo: ricerca/06_indice-ragionato.md, capitolo 14
% ------------------------------------------------------------------
\chapter{Il metodo in una settimana}
\label{cap:metodo}

\epigrafe[la memoria non deriva soltanto dalla natura, ma dipende moltissimo dall'arte e dall'impegno]{memoria non solum a natura proficiscitur, sed etiam habet plurimum artis et industriae}{Tommaso d'Aquino, Summa theologiae II-II, q. 49, a. 1, ad 2}

\iniziale{Q}{uesto capitolo} \segnaposto{apertura: da scrivere — vedi la scheda nell'indice ragionato}.

\section{Prima, durante e subito dopo}

\segnaposto{da scrivere}

\section{Il giorno dopo}

\segnaposto{da scrivere}

\section{Le settimane successive}

\segnaposto{da scrivere}

\section{Prima dell'esame}

\segnaposto{da scrivere}

\section{Una settimana tipo all'università}

\segnaposto{da scrivere}

\begin{daricordare}
\segnaposto{il principio del capitolo in due righe}
\end{daricordare}

\begin{domande}
  \item \segnaposto{domanda di recupero su questo capitolo}
  \item \segnaposto{domanda di applicazione a una materia che studi}
  \item \segnaposto{domanda di ripresa su un capitolo precedente}
\end{domande}

\begin{sintesi}
  \item \segnaposto{punto di sintesi}
\end{sintesi}

\finecapitolo
